%% kapa: exam=1 round=30 year=2019 period=1 form=A subject=real_estate subject_ko=부동산학원론 seq=5 number=85
%% source: Q-Net 감정평가사 기출문제 (cst00309) artlSeq=5205205 file=extracted/1008/5205205/2019년 제30회 감정평가사 1차시험 문제지/2019년 제30회 감정평가사 1차시험 문제지/2019년도 제30회 감정평가사 1차 1교시 A형.hwp
%% conversion: kordoc 4.24.51 -> tools/gen1.py (auto); review=proofread-sonnet; flags=none
\begin{question}{85}
토지정책에 관한 설명으로 옳은 것은?
\begin{choices}
  \item 토지정책수단 중 토지비축제도, 토지수용, 금융지원, 보조금 지급은 간접개입방식이다.
  \item 개발부담금제는 개발이 제한되는 지역의 토지소유권에서 개발권을 분리하여 개발이 필요한 다른 지역에 개발권을 양도할 수 있도록 하는 제도이다.
  \item 토지선매에 있어 시장ㆍ군수ㆍ구청장은 토지거래계약허가를 받아 취득한 토지를 그 이용목적대로 이용하고 있지 아니한 토지에 대해서 선매자에게 강제로 수용하게 할 수 있다.
  \item 개발권양도제는 개발사업의 시행으로 이익을 얻은 사업시행자로부터 개발이익의 일정액을 환수하는 제도이다.
  \item 토지적성평가제는 토지에 대한 개발과 보전의 경합이 발생했을 때 이를 합리적으로 조정하는 수단이다.
\end{choices}
\end{question}
\answer{5}
